\documentclass{article}

% A synthetic seed paper for Lattice's research environment. Its numbers are
% illustrative, not measured results.
\usepackage{iclr2026_conference,times}
\usepackage[utf8]{inputenc}
\usepackage[T1]{fontenc}
\usepackage{amsmath,amssymb,amsfonts}
\usepackage{mathtools}
\usepackage{bm}
\usepackage{booktabs}
\usepackage{tabularx}
\usepackage{graphicx}
\usepackage{hyperref}
\usepackage{url}
\usepackage{xcolor}

\newcommand{\loss}{\mathcal{L}}
\newcommand{\vect}[1]{\bm{#1}}
\newcommand{\todo}[1]{\textcolor{red}{#1}}

\title{Curriculum Signals in Self-Supervised Vision}
\author{
Seed Author \\
Lattice Research Environment \\
\texttt{seed@lattice.example}
}

\begin{document}

\maketitle

\begin{abstract}
Self-supervised vision models see their training images in random order.
We ask whether ordering images from easy to hard, with difficulty measured by the model's own contrastive loss, speeds up representation learning.
On a small benchmark, a loss-based curriculum reaches the final linear-probe accuracy of random ordering in fewer epochs, while hand-designed difficulty scores do not help.
\end{abstract}

\section{Introduction}
\label{sec:introduction}

Curriculum learning presents training examples in a meaningful order, typically from easy to hard \citep{bengio2009curriculum}.
Contrastive self-supervised methods such as SimCLR \citep{chen2020simple} and MoCo \citep{he2020momentum} learn by pulling together two augmented views of an image while pushing apart views of different images, and they usually sample images uniformly.
Section~\ref{sec:method} defines a curriculum driven by the contrastive loss itself, Section~\ref{sec:experiments} evaluates it, and Section~\ref{sec:discussion} discusses when ordering helps.

\section{Loss-based curriculum}
\label{sec:method}

Let $f_\theta$ be an encoder and $\vect{z}_i = f_\theta(\vect{x}_i)$ the representation of image $\vect{x}_i$.
For two views $\vect{z}_i, \vect{z}_i'$ of the same image and a batch of $N$ images, the contrastive loss is
\begin{equation}
  \loss_i \coloneqq -\log \frac{\exp(\vect{z}_i^\top \vect{z}_i' / \tau)}{\sum_{j=1}^{N} \exp(\vect{z}_i^\top \vect{z}_j' / \tau)},
  \label{eq:contrastive}
\end{equation}
with temperature $\tau$.
We use the running mean of $\loss_i$ over the last epoch as the difficulty $d_i$ of image $i$.

\paragraph{Schedule.}
At epoch $e$ of $E$, the sampler draws from the easiest fraction $\lambda(e)$ of images, where
\begin{equation}
  \lambda(e) = \min\!\left(1, \lambda_0 + (1 - \lambda_0)\frac{e}{E_\mathrm{warm}}\right),
  \label{eq:pacing}
\end{equation}
so that every image is eligible after $E_\mathrm{warm}$ epochs.
Difficulties are recomputed once per epoch, which adds one forward pass per image (Equation~\ref{eq:contrastive} needs no labels).
Table~\ref{tab:scores} lists the difficulty scores we compare.

\begin{table}[t]
  \centering
  \caption{Difficulty scores compared in Section~\ref{sec:experiments}.}
  \label{tab:scores}
  \begin{tabularx}{\linewidth}{@{}lX@{}}
    \toprule
    Score & Definition \\
    \midrule
    Random & no ordering; the standard baseline \\
    Entropy & Shannon entropy of the grey-level histogram, a hand-designed proxy \\
    Loss & running mean of $\loss_i$ from Equation~\ref{eq:contrastive} \\
    \bottomrule
  \end{tabularx}
\end{table}

\section{Experiments}
\label{sec:experiments}

We pretrain a ResNet-18 encoder \citep{he2016deep} for 200 epochs on a 100-class subset of images and evaluate it with a linear probe.\footnote{All numbers in this seed paper are illustrative and were not measured. The pacing parameters are $\lambda_0 = 0.3$ and $E_\mathrm{warm} = 60$, see Equation~\ref{eq:pacing}.}
Every configuration is run with three seeds.

\begin{figure}[t]
  \centering
  \includegraphics[width=0.9\linewidth]{figures/curriculum.png}
  \caption{Linear-probe accuracy during pretraining for random, entropy and loss ordering.}
  \label{fig:curriculum}
\end{figure}

Figure~\ref{fig:curriculum} shows that loss ordering reaches the random baseline's final accuracy after about 140 epochs.
Entropy ordering is slower than random throughout, which suggests that a poor difficulty signal is worse than none.
Table~\ref{tab:probe} summarises the final accuracies.

\begin{table}[t]
  \centering
  \caption{Final linear-probe accuracy (\%), mean $\pm$ standard deviation over three seeds (illustrative).}
  \label{tab:probe}
  \begin{tabular}{@{}lcc@{}}
    \toprule
    Ordering & 100 epochs & 200 epochs \\
    \midrule
    Random & $61.4 \pm 0.3$ & $68.2 \pm 0.2$ \\
    Entropy & $59.8 \pm 0.5$ & $67.5 \pm 0.4$ \\
    Loss & $\mathbf{64.0} \pm 0.4$ & $\mathbf{68.9} \pm 0.3$ \\
    \bottomrule
  \end{tabular}
\end{table}


\section{Discussion}
\label{sec:discussion}

The loss-based curriculum helps most early in training, when the model's notion of difficulty is still changing quickly (see the curve on page~\pageref{fig:curriculum}).
Hand-designed scores such as image entropy do not track what the model finds hard, which is consistent with the mixed results reported for fixed curricula in supervised learning \citep{wu2021curricula}.
A limitation of this study is its scale: we did not test whether the effect survives the much longer schedules used in practice.

\bibliography{references}
\bibliographystyle{iclr2026_conference}

\end{document}
